\section{Introduction}\label{sec:intro}

Emerging-market sovereign analysis sits awkwardly between three communities: macro-econometricians who chase causal identification, fixed-income strategists who price credit off a small set of fundamentals, and asset managers who want a tradable signal. The communities use different data, different time horizons, and different success metrics, and they rarely speak to each other in the same document. This paper benchmarks the development trajectory of 13 emerging economies --- the five BRICS, three members of the Alliance of Sahel States (Burkina Faso, Mali, Niger), Singapore, three Emerging-Africa cases (Rwanda, Kenya, Morocco), and Pakistan --- and runs that benchmark through all three lenses.

The pipeline has one anchor: a composite score that aggregates 39 development indicators across six sectors (education, energy, research \& innovation, health, housing \& living, security \& stability) into one number per country and year. The construction is principal-component within each sector and equal-weighted across sectors, on a panel that has been put through a per-capita and logarithmic transform to neutralise the country-size effect. Without that transform, sheer scale dominates: the raw composite places China at $+2.36$ and Singapore at $+2.04$ in 2024, because absolute patent counts and absolute scientific articles vary by three orders of magnitude. After transformation, the same composite places Singapore at $+2.93$ and China at $+2.33$, with Poland ($+2.29$) and Russia ($+1.98$) close behind. The new ranking is consistent with the consensus reading of where each economy sits today.

Three questions then follow.

\emph{First}, do KPI improvements cause macro outcomes? We estimate Jord{\`a} (2005) local projections, using the Hamilton (2018) cyclical component of each sector score as the shock and a two-way fixed-effect specification with Driscoll-Kraay standard errors. The IRFs at the chosen sample (23 countries, 35 years) are statistically weak: at horizon $h=1$, the cyclical-shock coefficients on log GDP per capita range from $-9.6$ (health) to $+6.7$ (security-stability), and the Driscoll-Kraay bands cover zero for every sector at most horizons. The Dumitrescu-Hurlin panel-Granger test, by averaging country-level Wald statistics rather than imposing a homogeneous $\beta_h$, reaches stronger conclusions: energy shocks Granger-cause poverty reduction ($p = 0.027$) and health shocks Granger-cause Gini reduction ($p = 0.046$), while no sector Granger-causes log GDP per capita at conventional levels. A synthetic-control estimate of Rwanda's Vision 2020 (treatment year 2000) yields a 2024 GDP-per-capita gap of $+\$1\,051$ PPP USD against a Burkina-Faso--Niger weighted counterfactual.

\emph{Second}, are KPIs already in the sovereign credit price? We compile the year-end S\&P foreign-currency long-term rating for 20 countries from 1995 to 2024 (the three Sahel members of AES and four others are unrated) and convert it to the standard 22-notch numeric scale. An expanding-window pooled ordered probit of the rating on the sector scores --- with no macro controls --- reaches an out-of-sample weighted MAE of 1.84 notches and an AUC of 0.81 for the investment-grade binary. The macro-only baseline (debt/GDP, current account, inflation, log GDPpc) reaches AUC 0.73. The full model is no better than the sector-only model in AUC and only marginally better in MAE. A Bayesian hierarchical ordered-logistic counterpart converges cleanly (NUTS, 4 chains $\times$ 1\,500 draws) and localises the signal to two sectors with credibly non-zero population loadings: energy ($\mu_\beta = +0.92$, 90~\% HDI $[+0.10,\,+1.74]$) and health ($\mu_\beta = -0.92$, HDI $[-1.38,\,-0.45]$); the four other sectors have HDIs that cover zero. The reading is direct: development KPIs subsume a sizeable share of the information that drives sovereign credit; macro controls add little once KPIs are in.

\emph{Third}, do KPIs generate tradable equity returns? We sort 17 country ETFs each month-end by their lagged composite score (six-month publication lag to rule out look-ahead) and form a long-only top-quintile portfolio with 20~bps per-side transaction costs. From January 2007 to December 2024 the strategy returns 3.6~\% annualised at 23.8~\% volatility (Sharpe 0.15, maximum drawdown 56~\%); the MSCI Emerging Markets benchmark (EEM) over the same window returns 2.6~\% at Sharpe 0.12. A non-parametric benchmark of 5\,000 random four-country baskets has mean Sharpe 0.28, placing the KPI strategy at the $0^{\text{th}}$ percentile of the random distribution --- below every random basket. A long-short top-minus-bottom variant returns $-11.1$~\% annualised, with most of the damage coming from a single dislocation in 2015 and the Russian delisting in early 2022. KPIs actively mis-allocate in equity space at this horizon and with this universe.

These three answers fit together into one observation. KPIs are slow-moving and so push their predictive power into instruments that price slowly: ratings, not weekly equity returns. The paper is structured to make that point.

Section~\ref{sec:literature} places the work in the credit-fundamentals, slow-moving-factor and panel-LP literatures. Section~\ref{sec:data} documents the panel, the sources, the imputation choices and the coverage matrix. Section~\ref{sec:methodology} presents the composite-scoring, LP, credit and portfolio specifications in one place. Sections~\ref{sec:partI}--\ref{sec:partIII} deliver the three blocks of results. Section~\ref{sec:robustness} reports rank-stability under jackknife, sensitivity to imputation, sub-period splits and a random-portfolio benchmark. Section~\ref{sec:conclusion} closes.

The accompanying repository (\texttt{newperformers}) reproduces every number in this paper from a single command (\verb|make all|). Where the text cites a coefficient, an AUC or a Sharpe ratio, the underlying CSV is in \verb|outputs/tables/| and the figure in \verb|outputs/figures/|.
